\documentclass[aspectratio=169, lualatex, handout]{beamer}
\makeatletter
\def\input@path{{theme/}}
\makeatother
\usetheme{cipher}

\title{Introduction to Cryptography}
\author{Nadim Kobeissi}
\institute{American University of Beirut \\ \url{https://appliedcryptography.page}}
\date{\today}

\coverpartname{Part 1: Cryptography in Theory}
\coversessionname{1.3: Rudiments of Provable Security}

\begin{document}
	\begin{frame}[plain]
		\titlepage
	\end{frame}

	\begin{frame}{Last time, we defined subroutines}
		\begin{center}
			\begin{tikzpicture}[
				box/.style={rectangle, draw, minimum height=2.5cm, align=left, fill=white},
				dashed box/.style={rectangle, draw, dashed, minimum width=2.5cm, minimum height=2.5cm, align=center, fill=white},
				>=Stealth,
			]
				\node[dashed box] (adversary) {$adversary$};
				\node[box, right=2cm of adversary](attack){
					$\underline{\textsc{Attack}\scalebox{1.1}{(M)}}$: \hfill \quad \quad \textcolor{gray}{\scriptsize \textit{// adversary chooses $M$}} \\
					\quad $K \twoheadleftarrow \bits^{n}$ \hfill \quad \quad \textcolor{gray}{\scriptsize \textit{// victim samples $K$}} \\
					\quad $C \coloneq \textsf{Enc}(K, M)$ \hfill \quad \quad \textcolor{gray}{\scriptsize \textit{// victim encrypts}} \\
					\quad return $C$ \hfill \quad \quad \textcolor{gray}{\scriptsize \textit{// adversary sees $C$}}
				};
				\draw[->] (adversary) -- node[above] {$M$} (attack);
				\draw[<-] ([yshift=-0.5cm]adversary.east) -- node[below] {$C$} ([yshift=-0.5cm]attack.west);
			\end{tikzpicture}
		\end{center}
	\end{frame}

	\begin{frame}{Subroutines}
		\begin{columns}[c]
			\begin{column}{0.5\textwidth}
				\begin{itemize}[<+->]
					\item ``Victim'' chooses their key.
					\item Adversary chooses the message and receives the ciphertext.
					\item We say that \textbf{the adversary has access to an encryption oracle}.
				\end{itemize}
			\end{column}
			\begin{column}{0.5\textwidth}
				\imagewithcaption{attacker_interface.svg}{Source: The Joy of Cryptography}
			\end{column}
		\end{columns}
	\end{frame}

	\begin{frame}{Attack scenarios as libraries}
		\begin{columns}[t]
			\begin{column}{0.65\textwidth}
				\begin{itemize}[<+->]
					\item This is a \textbf{library} with two \textbf{subroutines} and a \textbf{global variable} $D$.
					\item Victim holds a 6-sided die.
					\item The first line of the library represents an initialization step.
					\item Attacker can call the subroutines at any time, and:
				 	\begin{enumerate}
						\item Make a guess about the current value of the die and learn whether the guess was correct.
						\item Instruct the victim to (privately) re-roll the die.
					\end{enumerate}
				\end{itemize}
			\end{column}
			\begin{column}{0.35\textwidth}
				\vspace{-1cm}
				\begin{center}
					\sslibrary{dice-guess}{
						$D \twoheadleftarrow \{1, 2, 3, 4, 5, 6\}$ \\[1em]
						\sslibrarysubroutine{Guess}{G}{
							return $D \equiv G$
						}{1} \\[1em]
						\sslibrarysubroutine{Reroll}{G}{
							$D \twoheadleftarrow \{1, 2, 3, 4, 5, 6\}$
						}{1}
					}
				\end{center}
			\end{column}
		\end{columns}
	\end{frame}

	\begin{frame}{Test}
		\begin{columns}[c]
			\begin{column}{0.5\textwidth}
				\begin{itemize}
					\item Test
				\end{itemize}
			\end{column}
			\begin{column}{0.5\textwidth}
				\sslinked{
					\ssprogram{}{
						if \sslibraryfun{Guess}{6}:\\
						\quad return true\\
						\sslibraryfun{Reroll}{}\\
						return \sslibraryfun{Guess}{6}
					}{0.75}
				}{
					\sslibrary{dice-guess}{
						$D \twoheadleftarrow \{1, 2, 3, 4, 5, 6\}$ \\[1em]
						\sslibrarysubroutine{Guess}{G}{
							return $D \equiv G$
						}{1} \\[1em]
						\sslibrarysubroutine{Reroll}{G}{
							$D \twoheadleftarrow \{1, 2, 3, 4, 5, 6\}$
						}{0.75}
					}
				}
			\end{column}
		\end{columns}
	\end{frame}

	\begin{frame}[plain]
		\titlepage
	\end{frame}
\end{document}
